\documentclass[12pt]{article}
\usepackage[margin=0.75in]{geometry}
\usepackage{fontspec}
\setmainfont{Latin Modern Roman}
\usepackage{microtype}
\usepackage{multicol}
\usepackage{fancyhdr}
\usepackage{titlesec}
\usepackage{enumitem}
\usepackage{xcolor}
\usepackage[hidelinks]{hyperref}
\setlength{\columnsep}{0.28in}
\setlength{\parindent}{1.1em}
\setlength{\parskip}{0.32em}
\setlength{\headheight}{15pt}
\titleformat{\section}{\large\bfseries\scshape}{}{0pt}{}
\titleformat{\subsection}{\normalsize\bfseries}{}{0pt}{}
\pagestyle{fancy}
\fancyhf{}
\fancyfoot[C]{\thepage}
\renewcommand{\headrulewidth}{0.4pt}
\newcommand{\focusbox}[1]{\par\vspace{0.4em}\noindent\begingroup\setlength{\fboxsep}{6pt}\colorbox{gray!12}{\begin{minipage}{\dimexpr\linewidth-2\fboxsep\relax}#1\end{minipage}}\endgroup\par\vspace{0.5em}}

\fancyhead[L]{Whately: Fallacies}
\fancyhead[R]{Student Reading}
\begin{document}
\begin{center}
{\LARGE\bfseries What Is a Fallacy?}\par\vspace{0.25em}
{\large\scshape Week 7: When an Argument Looks Stronger Than It Is}\par\vspace{0.5em}
{Adapted and modernized from Richard Whately, \textit{Elements of Logic}}\par\vspace{0.6em}
\rule{0.82\textwidth}{0.4pt}
\end{center}

\section*{Central Question}
Why do bad arguments often feel persuasive?

\section*{Vocabulary Preview}
\begin{multicols}{2}
\begin{itemize}[leftmargin=*, itemsep=0.18em]
\item \textbf{Fallacy}: reasoning that seems to prove a conclusion but fails to prove it.
\item \textbf{Appearance of proof}: the surface look of support that can persuade without earning the conclusion.
\item \textbf{Defect in reasoning}: the real break in the argument---where the support does not reach the claim.
\item \textbf{False claim}: a statement that is not true.
\item \textbf{Bad inference}: a move from reasons to conclusion that does not follow, even if some claims are true.
\item \textbf{Emotional force}: the power of feeling---fear, shame, anger, loyalty---to make a claim feel settled.
\item \textbf{Rhetorical force}: the power of skillful speaking, repetition, tone, and crowd pressure to make a claim feel settled.
\item \textbf{Logical force}: the power of actual support: whether the reasons, if granted, justify the conclusion.
\item \textbf{Fallacy autopsy}: a disciplined diagnosis that separates what the argument seems to prove from where it fails.
\item \textbf{Persuasive nonsense}: language that moves an audience without establishing a conclusion.
\item \textbf{Question at issue}: the exact claim that needed proof in this debate.
\item \textbf{Repair}: rewriting the argument with a fairer claim or clearer support.
\end{itemize}
\end{multicols}

\section*{Introduction}
In Weeks 1--6 you learned to parse an argument, fix unstable terms, state clear propositions, map a syllogism, test formal validity, and separate form from fact. You can now ask: What is the conclusion? What reasons are offered? Do the terms stay steady? Does the form work? Are the premises true?

This week opens a new chapter in Whately's design: \textit{fallacies}. A fallacy is not merely a false sentence. It is reasoning that \textit{seems} to prove something while failing to prove it. Bad arguments often feel strong because they borrow the appearance of proof from emotion, confidence, reputation, clever wording, or crowd pressure. The disciplined thinker's job is to slow the argument down and name the defect.

\focusbox{\textbf{Source note:} This reading is adapted and modernized from Whately's opening treatment of fallacies in \textit{Elements of Logic}. It preserves his concern with the \textit{appearance} of proof and the difference between persuasion and real support. Classroom examples are original adaptations for the twelfth-grade long logic unit. Later weeks will name more specific fallacy families; this week teaches the general habit of diagnosis.}

\begin{multicols}{2}
\section*{Reading}

\subsection*{1. Not Every Mistake Is a Fallacy}
A person can be wrong without committing a fallacy. If someone says, ``The bus arrives at 7:10,'' and the bus arrives at 7:25, the claim is false. That is a factual error. It is not yet a fallacy, because no chain of reasoning has been offered that only \textit{looks} like proof.

A fallacy appears when someone offers reasons that seem to support a conclusion but do not actually support it. The danger is not only that the speaker is wrong. The danger is that the reasoning \textit{feels} finished when it is not.

\subsection*{2. The Appearance of Proof}
Whately's key insight is simple and hard: fallacies work because they look better than they are. An argument may sound confident, moral, patriotic, scientific, witty, or kind. It may use the vocabulary of proof---``therefore,'' ``obviously,'' ``everyone knows,'' ``no true friend would deny.'' It may still fail to establish the conclusion.

The first question is not, ``Do I like this speaker?'' The first question is, ``What is being proved, and has it been proved?''

\subsection*{3. False Claim Versus Bad Inference}
Two failures must be kept separate.

\begin{enumerate}[leftmargin=*,itemsep=0.12em]
\item \textbf{False claim:} one of the statements is untrue.
\item \textbf{Bad inference:} the conclusion does not follow from the reasons given.
\end{enumerate}

Example of a mainly false claim: ``All transfer students skip class. Maya is a transfer student. Therefore Maya skips class.'' The form may look tidy, but the first premise is false. The argument fails on the fact side.

Example of a mainly bad inference: ``Jordan spoke against the new attendance policy. Therefore Jordan does not care about graduation.'' Even if Jordan did speak against the policy, the conclusion does not follow. Many careful seniors can oppose a policy for serious reasons. Here the defect is the leap, not merely a calendar fact.

A strong diagnosis names which kind of failure is doing the real damage---or whether both are present.

\subsection*{4. Emotional Force and Rhetorical Force}
Arguments live among people, not only on chalkboards. Shame, loyalty, fear, admiration, and anger can make a conclusion feel inevitable. Skillful speaking can do the same: repetition, irony, flattery, group pressure, and vivid images can make weak support feel strong.

Emotional force and rhetorical force are real. They can move a crowd. They are not the same as logical force. Logical force asks whether the reasons, if granted, justify the conclusion. A disciplined thinker can feel the pressure of a speech and still ask whether the proof has been earned.

\subsection*{5. Why Fallacies Feel Strong}
Fallacies often feel strong for ordinary human reasons:

\begin{itemize}[leftmargin=*,itemsep=0.12em]
\item they flatter what we already want to believe;
\item they attack a person instead of testing a claim;
\item they prove a nearby claim while missing the real question;
\item they hide a weak step under confident language;
\item they mix true statements with an unsupported leap;
\item they use a word in two senses without fixing the term.
\end{itemize}

You already have tools for some of these: the four-step test, term fixing, proposition clarity, syllogism maps, validity, and soundness. Fallacy work does not throw those tools away. It uses them when an argument is trying to \textit{look} like proof.

\subsection*{6. The Fallacy Autopsy}
When an argument seems persuasive but suspicious, perform a short autopsy:

\begin{enumerate}[leftmargin=*,itemsep=0.12em]
\item \textbf{Conclusion:} What is the speaker trying to make you accept?
\item \textbf{Apparent proof:} What reasons, tone, or pressure make the conclusion feel established?
\item \textbf{Defect:} Where does the reasoning fail---false claim, bad inference, unstable term, missing question, unproved pressure point, or overclaim?
\item \textbf{Repair or narrow:} What fairer claim could be made, or what evidence would actually be needed?
\end{enumerate}

This autopsy is the Week 7 habit. Later weeks will name particular families of fallacy more carefully. This week trains the eye to see \textit{appearance} versus \textit{support}.

\subsection*{7. A School Case}
Consider this argument in a student council meeting: ``Anyone who questions the spirit-week schedule is against school pride. You questioned the schedule. Therefore you are against school pride.''

The argument can feel strong in a crowded room. It wraps disagreement in moral accusation. But the first premise is doubtful, and the leap from questioning a schedule to opposing school pride is a bad inference. The conclusion may silence the speaker without proving the charge.

\subsection*{8. Macbeth and Apparent Proof}
In \textit{Macbeth}, persuasion often works by pressure rather than proof. Lady Macbeth does not carefully demonstrate that murder is justified. She attacks Macbeth's courage and manhood so that hesitation feels like shame. The speech can move Macbeth because it is fierce, intimate, and humiliating---not because it has earned the conclusion that Duncan must die.

The literary question for this week is not only what characters desire. It is how language creates the \textit{appearance} of necessity. Emotional force can settle a mind before logical force has done any work.

\subsection*{9. What This Week Adds}
Week 7 teaches you to diagnose arguments that look stronger than they are. A complete Week 7 diagnosis should include:

\begin{itemize}[leftmargin=*,itemsep=0.12em]
\item the conclusion;
\item what gives the argument its appearance of proof;
\item whether the main failure is a false claim, a bad inference, or both;
\item where emotional or rhetorical force is doing work that logic has not done;
\item a repaired claim or a statement of what evidence would actually be needed.
\end{itemize}

This habit protects you as a reader, a citizen, and a writer. It also prepares the next weeks, when specific fallacy types---equivocation in full arguments, circular reasoning, irrelevant proof, and more---will be named and practiced.

\section*{Reading Focus}
\begin{enumerate}[leftmargin=*]
\item In your own words, what is a fallacy?
\item Why is a false claim not automatically a fallacy?
\item What is the difference between emotional force and logical force?
\item Explain the difference between a false claim and a bad inference.
\item Why do fallacies often feel persuasive even after they are exposed?
\item List the four steps of a fallacy autopsy.
\item How does Lady Macbeth's pressure on Macbeth illustrate the appearance of proof?
\end{enumerate}

\section*{Handout Summary}
Write 5--7 sentences explaining why bad arguments can feel strong and how a disciplined thinker should respond. Your summary must include the words \textit{fallacy}, \textit{appearance}, \textit{inference}, and \textit{proof}.
\end{multicols}
\end{document}
